% Arti Lambang dan Singkatan -- transcribed verbatim from
% Daftar/Arti Lambang dan Singkatan.doc
\chapter*{ARTI LAMBANG DAN SINGKATAN}
\addcontentsline{toc}{chapter}{ARTI LAMBANG DAN SINGKATAN}

\begin{description}[style=nextline,leftmargin=2.2cm,labelwidth=2cm]
\item[$a$] Tinggi blok tegangan persegi ekivalen
\item[$AC$] Beban titik kumpul gabungan
\item[$AE$] Beban titik kumpul ekivalen
\item[$AFC$] Beban titik kumpul gabungan yang selaras dengan $DF$
\item[$AM_i$] Gaya di ujung batang $i$
\item[$AML$] Gaya ujung batang akibat beban
\item[$AMS_i$] Gaya jepit kedua ujung $i$ dalam arah sumbu struktur
\item[$AR$] Reaksi di titik kumpul yang dikekang
\item[$ARC$] Beban titik kumpul gabungan yang selaras dengan perpindahan pada titik kumpul yang dikekang
\item[$AS$] Luas tulangan pada daerah tarik
\item[$AS'$] Luas tulangan pada daerah tekan
\item[$b$] Lebar penampang beton
\item[$B$] Lebar kolom
\item[$c$] Jarak dari serat tekan terluar ke garis netral
\item[$CC$] Gaya internal pada daerah tekan beton
\item[$CS$] Gaya internal pada daerah tekan baja tulangan
\item[$d$] Jarak dari serat tekan terluar ke pusat tulangan tarik
\item[$d'$] Jarak dari serat tarik terluar ke pusat tulangan desak
\item[$dia$] Diameter tulangan pada kolom
\item[$diatldslp$] Diameter tulangan desak pada daerah lapangan balok
\item[$diatldstm$] Diameter tulangan desak pada daerah tumpuan balok
\item[$diatlgs$] Diameter tulangan geser
\item[$diatltrlp$] Diameter tulangan tarik pada daerah lapangan balok
\item[$diatltrtm$] Diameter tulangan tarik pada daerah tumpuan balok
\item[$DF$] Perpindahan titik kumpul bebas
\item[$DM_i$] Perpindahan di ujung batang $i$
\item[$DR$] Perpindahan titik kumpul yang dikekang
\item[$E$] Modulus elastisitas untuk gaya aksial
\item[$E_C$] Modulus elastisitas untuk beton
\item[$E_S$] Modulus elastisitas baja tulangan
\item[$f$] Fungsi sasaran
\item[$f_c'$] Kuat desak karakteristik beton
\item[$f_S'$] Tegangan yang terjadi pada baja tarik
\item[$f_y$] Kuat tarik baja tulangan
\item[$g_j$] Fungsi kendala pertaksamaan ke-$j$
\item[$G$] Modulus elastisitas untuk geser
\item[$h_j$] Fungsi kendala persamaan ke-$j$
\item[$H$] Tinggi kolom
\item[$I_x$] Momen inersia pada sumbu $x$
\item[$I_y$] Momen inersia pada sumbu $y$
\item[$I_z$] Momen inersia pada sumbu $z$
\item[$jaraktlgs$] Jarak tulangan geser
\item[$JVD$] Jumlah variabel desain
\item[$L$] Panjang bentang batang
\item[$m$] Jumlah batang
\item[$M$] Momen
\item[$M_n$] Momen tahanan nominal penampang
\item[$M_{ox}$] Momen desain efektif pada arah $x$
\item[$M_{oy}$] Momen desain efektif pada arah $y$
\item[$M_x$] Momen yang terjadi pada arah sumbu $x$
\item[$M_y$] Momen yang terjadi pada arah sumbu $y$
\item[$nB$] Jumlah balok pada struktur
\item[$nK$] Jumlah kolom pada struktur
\item[$ntldslp$] Jumlah tulangan desak pada daerah lapangan balok
\item[$ntldstm$] Jumlah tulangan desak pada daerah tumpuan balok
\item[$ntltrlp$] Jumlah tulangan tarik pada daerah lapangan balok
\item[$ntltrtm$] Jumlah tulangan tarik pada daerah tumpuan balok
\item[$N$] Jumlah tulangan pada salah satu sisi kolom
\item[$Ntitik$] Jumlah titik coba
\item[$P$] Gaya aksial
\item[$P_u$] Gaya aksial ultimit
\item[$SFF$] Gaya $AF$ akibat satu satuan perpindahan $DF$
\item[$SJ$] Matrik kekakuan titik terakit
\item[$SM$] Kekakuan batang dalam arah sumbu struktur
\item[$SMS_i$] Kekakuan batang untuk kedua ujung batang $i$ dalam arah sumbu struktur
\item[$SRF$] Reaksi $AR$ akibat satu perpindahan $DF$
\item[$T$] Gaya internal pada daerah tarik baja tulangan
\item[$x$] Variabel desain
\item[$X$] Ruang variabel desain
\item[$X_B$] Koordinat $x$ titik terbaik
\item[$X_G$] Koordinat $x$ titik \textit{good}
\item[$X_M$] Koordinat $x$ titik tengah
\item[$X_S$] Koordinat $x$ arah penelusuran
\item[$X_T$] Koordinat $x$ titik coba baru
\item[$X_W$] Koordinat $x$ titik terjelek
\item[$Y_B$] Koordinat $y$ titik terbaik
\item[$Y_G$] Koordinat $y$ titik \textit{good}
\item[$Y_M$] Koordinat $y$ titik tengah
\item[$Y_S$] Koordinat $y$ arah penelusuran
\item[$Y_T$] Koordinat $y$ titik coba baru
\item[$Y_W$] Koordinat $y$ titik terjelek
\item[$\alpha$] Koefisien biaksial
\item[$\beta_1$] Faktor reduksi tinggi blok tegangan
\item[$\varepsilon$] Faktor konvergensi
\item[$\varepsilon_s$] Regangan pada tulangan tarik
\item[$\varepsilon_s'$] Regangan pada tulangan desak
\item[$\phi$] Faktor reduksi kekuatan
\item[$\rho$] Rasio penulangan pada daerah tarik
\item[$\rho'$] Rasio penulangan pada daerah tekan
\item[$\rho_b$] Rasio imbang
\item[$\rho_{maks}$] Rasio penulangan maksimum yang diijinkan
\end{description}

\clearpage
